% ══════════════════════════════════════════════════════════════
% Chapter 14 -- Realized Covariance and Multivariate Forecasting
% ══════════════════════════════════════════════════════════════

\chapter{Realized Covariance and Multivariate Forecasting}\label{ch:multivariate}

\begin{application}[From One Asset to Many]
Chapters~\ref{ch:returns}--\ref{ch:vrp} focused on the volatility of a single asset.
In practice, portfolios hold many assets, and the \emph{covariance structure} matters as
much as individual volatilities. Minimum-variance portfolios, risk parity, hedging, and
option pricing on baskets all require a full covariance matrix, not just a list of
variances. This chapter extends realized volatility to the multivariate setting:
estimating, modeling, and forecasting entire covariance matrices. Your
project uses this material directly: the cross-asset covariance structure
feeds the Layer~4 spillover features that enter LightGBM as scalar inputs.
\end{application}

% ──────────────────────────────────────────────────────────────
\section{Realized Covariance}\label{sec:multivariate:rc}
% ──────────────────────────────────────────────────────────────

\textbf{Where we are.}  You know how to build realized variance for one asset
(Chapter~\ref{ch:rv}). The natural multivariate extension replaces squared returns
with outer products.

\begin{definition}[Realized Covariance Matrix]
Let $\mathbf{r}_{t,i} \in \R^p$ be the $i$-th intraday return vector on day $t$,
for $i = 1, \ldots, M$. The \textbf{realized covariance} (RC) matrix is
\begin{equation}\label{eq:rc}
  \mathbf{RC}_t \;=\; \sum_{i=1}^{M} \mathbf{r}_{t,i}\,\mathbf{r}_{t,i}^\top
  \;\in\; \R^{p \times p}.
\end{equation}
\begin{itemize}[nosep]
  \item $\mathbf{r}_{t,i} \in \R^p$: intraday return vector for $p$ assets at
        interval $i$.
  \item $M$: number of intraday intervals (e.g., 78 for 5-minute sampling of
        a 6.5-hour trading day).
  \item Diagonal entries $[\mathbf{RC}_t]_{jj}$: realized variances of each asset
        (as in Chapter~\ref{ch:rv}).
  \item Off-diagonal entries $[\mathbf{RC}_t]_{jk}$: realized covariances between
        assets $j$ and $k$.
\end{itemize}
Under no noise and synchronous trading, $\mathbf{RC}_t$ converges to the
integrated covariance matrix as $M \to \infty$
\citep{BarndorffNielsenHansenLundeShephard2011}.
\end{definition}

% ── TikZ: Covariance matrix heatmap schematic ──
\begin{figure}[ht]
\centering
\begin{tikzpicture}[scale=0.45]
  % Draw the matrix grid
  \def\n{10}

  % Sector blocks - Tech (assets 1-4), Finance (5-7), Energy (8-10)
  % High covariance (dark fill) within sector, low (light fill) across sector

  % Tech block (top-left 4x4) - strong covariance
  \foreach \i in {0,...,3} {
    \foreach \j in {0,...,3} {
      \pgfmathsetmacro{\shade}{ifthenelse(\i==\j, 90, 55+5*rnd)}
      \fill[defblue!\shade!white] (\i, \n-1-\j) rectangle ++(1,1);
    }
  }

  % Finance block (middle 3x3) - strong covariance
  \foreach \i in {4,...,6} {
    \foreach \j in {4,...,6} {
      \pgfmathsetmacro{\shade}{ifthenelse(\i==\j, 90, 50+5*rnd)}
      \fill[keyorange!\shade!white] (\i, \n-1-\j) rectangle ++(1,1);
    }
  }

  % Energy block (bottom-right 3x3) - strong covariance
  \foreach \i in {7,...,9} {
    \foreach \j in {7,...,9} {
      \pgfmathsetmacro{\shade}{ifthenelse(\i==\j, 90, 55+5*rnd)}
      \fill[warnred!\shade!white] (\i, \n-1-\j) rectangle ++(1,1);
    }
  }

  % Cross-sector blocks - weak covariance (light gray)
  \foreach \i in {0,...,3} {
    \foreach \j in {4,...,6} {
      \fill[gray!15] (\i, \n-1-\j) rectangle ++(1,1);
      \fill[gray!15] (\j, \n-1-\i) rectangle ++(1,1);
    }
  }
  \foreach \i in {0,...,3} {
    \foreach \j in {7,...,9} {
      \fill[gray!15] (\i, \n-1-\j) rectangle ++(1,1);
      \fill[gray!15] (\j, \n-1-\i) rectangle ++(1,1);
    }
  }
  \foreach \i in {4,...,6} {
    \foreach \j in {7,...,9} {
      \fill[gray!20] (\i, \n-1-\j) rectangle ++(1,1);
      \fill[gray!20] (\j, \n-1-\i) rectangle ++(1,1);
    }
  }

  % Grid lines
  \draw[gray!50, thin] (0,0) grid (\n, \n);
  \draw[thick] (0,0) rectangle (\n, \n);

  % Sector labels
  \draw[decorate, decoration={brace, amplitude=5pt, mirror}]
    (0, -0.3) -- (4, -0.3) node[midway, below=6pt, font=\small]{Tech};
  \draw[decorate, decoration={brace, amplitude=5pt, mirror}]
    (4, -0.3) -- (7, -0.3) node[midway, below=6pt, font=\small]{Finance};
  \draw[decorate, decoration={brace, amplitude=5pt, mirror}]
    (7, -0.3) -- (10, -0.3) node[midway, below=6pt, font=\small]{Energy};

  % Annotations
  \node[right, font=\small, align=left] at (11, 8) {Dark = high covariance};
  \node[right, font=\small, align=left] at (11, 7) {Light = low covariance};
  \node[right, font=\small, align=left] at (11, 5.5) {Block structure:};
  \node[right, font=\small, align=left] at (11, 4.7) {within-sector pairs};
  \node[right, font=\small, align=left] at (11, 3.9) {co-move strongly};

  % Dimension annotation
  \draw[<->, thick] (0, 10.5) -- (10, 10.5) node[midway, above, font=\small]{$p$ assets};
\end{tikzpicture}
\caption{Schematic of a realized covariance matrix with $p = 10$ assets sorted
by sector. The block-diagonal structure (strong within-sector covariance, weak
cross-sector) is a key pattern that CNNs and factor models exploit. The diagonal
entries are the realized variances from Chapter~\ref{ch:rv}.}
\label{fig:rc-heatmap}
\end{figure}

\subsection{The Synchronicity Problem}

Forming each $\mathbf{r}_{t,i}$ requires observing all $p$ assets at the same
timestamps, but different assets trade at different times. A
thinly traded stock might have no trade in a given 5-minute window. If you
simply ignore missing observations, you bias realized covariance toward zero
(the ``Epps effect'').

% ── TikZ: Non-synchronous trading ──
\begin{figure}[ht]
\centering
\begin{tikzpicture}[>=Stealth, xscale=1.1]
  % time axis
  \draw[->, thick] (0,0) -- (11.5,0) node[right]{\small time};

  % Asset A trades
  \node[left] at (0,1.2) {\textbf{Asset A}};
  \foreach \x in {0.5, 2.0, 3.8, 5.5, 7.0, 8.2, 10.0, 11.0} {
    \draw[defblue, thick] (\x, 0.9) -- (\x, 1.5);
    \fill[defblue] (\x, 1.2) circle (2pt);
  }

  % Asset B trades
  \node[left] at (0,2.4) {\textbf{Asset B}};
  \foreach \x in {1.0, 1.8, 4.2, 6.0, 6.8, 9.5, 10.5} {
    \draw[warnred, thick] (\x, 2.1) -- (\x, 2.7);
    \fill[warnred] (\x, 2.4) circle (2pt);
  }

  % Refresh times
  \foreach \x in {1.0, 2.0, 4.2, 6.0, 7.0, 9.5, 10.5} {
    \draw[intgreen, dashed] (\x, -0.3) -- (\x, 3.0);
  }
  \node[intgreen, below] at (5.5, -0.5) {\small Refresh times (both assets have traded)};
\end{tikzpicture}
\caption{Non-synchronous trading. Asset A (blue) and Asset B (red) trade at
different times. Refresh times (green dashed) mark points where both assets
have at least one new observation since the last refresh time.}
\label{fig:nonsync}
\end{figure}

\subsection{Refresh-Time Sampling}

The simplest fix: wait until \emph{all} $p$ assets have traded at least once,
then record a synchronized return vector. These ``refresh times'' are the green
lines in Figure~\ref{fig:nonsync}. Formally, define $\tau_0 = 0$ and
\[
  \tau_k \;=\; \max_{j=1,\ldots,p}\; \min\bigl\{t_{j,n} : t_{j,n} > \tau_{k-1}\bigr\},
\]
where $t_{j,n}$ is the $n$-th trade time for asset $j$. You compute returns at
refresh times and plug them into \eqref{eq:rc}. The cost: you lose observations,
sometimes dramatically for illiquid assets.

\subsection{Hayashi--Yoshida Estimator}

\citet{HayashiYoshida2005} proposed a more efficient solution. Instead of aligning
timestamps, you sum the products of all \emph{overlapping} returns.

\begin{definition}[Hayashi--Yoshida Estimator]
For two assets with (possibly different) trade times, let $r_{t,i}^{(A)}$ denote
asset A's return over interval $[t_{i-1}^A, t_i^A)$ and similarly for B. The
Hayashi--Yoshida realized covariance is
\begin{equation}\label{eq:hy}
  \widehat{\sigma}_{AB,t}^{\mathrm{HY}}
  \;=\; \sum_{i}\sum_{j}\; r_{t,i}^{(A)}\, r_{t,j}^{(B)}\;
  \mathbf{1}\!\Bigl\{
    [t_{i-1}^A, t_i^A) \cap [t_{j-1}^B, t_j^B) \neq \emptyset
  \Bigr\}.
\end{equation}
\begin{itemize}[nosep]
  \item Sum over all pairs of intervals that overlap in time.
  \item No data is discarded; every trade contributes.
  \item Unbiased and consistent for the integrated covariance under mild conditions.
  \item Not guaranteed positive semi-definite (PSD) for $p > 2$.
\end{itemize}
\end{definition}

\begin{intuition}[In Plain English]
Instead of forcing both assets onto a common clock (which throws away data),
Hayashi--Yoshida is like computing a dot product, but one that uses every
available trade rather than only the synchronized ones.
\end{intuition}

\begin{workedexample}{Hayashi--Yoshida with 2 Assets}
Suppose on one day, asset A trades at times $\{0, 2, 5, 8\}$ with prices
$\{100, 101, 99, 100.5\}$ and asset B trades at times $\{0, 3, 6, 8\}$ with
prices $\{50, 50.5, 49.8, 50.2\}$.

\textbf{Step 1: compute returns.}
\begin{align*}
  r_1^A &= \ln(101/100) \approx 0.00995, &\quad
  r_2^A &= \ln(99/101) \approx -0.02000, &\quad
  r_3^A &= \ln(100.5/99) \approx 0.01504, \\
  r_1^B &= \ln(50.5/50) \approx 0.00995, &\quad
  r_2^B &= \ln(49.8/50.5) \approx -0.01396, &\quad
  r_3^B &= \ln(50.2/49.8) \approx 0.00800.
\end{align*}

\textbf{Step 2: identify overlapping pairs.}
\begin{center}
\begin{tabular}{ccc}
\toprule
$r_i^A$ interval & $r_j^B$ interval & Overlap? \\
\midrule
$[0,2)$ & $[0,3)$ & Yes \\
$[0,2)$ & $[3,6)$ & No  \\
$[2,5)$ & $[0,3)$ & Yes \\
$[2,5)$ & $[3,6)$ & Yes \\
$[5,8)$ & $[3,6)$ & Yes \\
$[5,8)$ & $[6,8)$ & Yes \\
\bottomrule
\end{tabular}
\end{center}

\textbf{Step 3: sum products of overlapping pairs.}
\[
  \widehat{\sigma}_{AB}^{\mathrm{HY}} \approx 9.0 \times 10^{-5}.
\]
The positive value indicates slight positive comovement. Note: no observations
were discarded.
\end{workedexample}

\begin{warning}[PSD is Not Guaranteed]
For $p = 2$, the Hayashi--Yoshida estimator is always PSD (it is a scalar
covariance). For $p \geq 3$, the matrix assembled from pairwise HY estimates
can have negative eigenvalues. You may need to project the result onto the
PSD cone (e.g., by zeroing out negative eigenvalues), which introduces bias.
Section~\ref{sec:multivariate:psd} discusses this systematically.
\end{warning}


% ──────────────────────────────────────────────────────────────
\section{Multivariate Realized Kernel}\label{sec:multivariate:mrk}
% ──────────────────────────────────────────────────────────────

\textbf{Where we are.}  You have two problems: non-synchronous trading
(Section~\ref{sec:multivariate:rc}) and microstructure noise
(Chapter~\ref{ch:noise}). The multivariate realized kernel handles both
simultaneously.

\begin{intuition}[Extending the Realized Kernel]
In Chapter~\ref{ch:noise}, the univariate realized kernel applied a weighting
function to autocovariances of returns at different lags, suppressing noise
while remaining consistent. The multivariate version does exactly the same
thing, but with cross-autocovariance matrices instead of scalar autocovariances.
\end{intuition}

\begin{definition}[Multivariate Realized Kernel]
\citet{BarndorffNielsenHansenLundeShephard2011} define the multivariate
realized kernel as
\begin{equation}\label{eq:mrk}
  \mathbf{K}_t \;=\; \sum_{h=-H}^{H} k\!\left(\frac{h}{H+1}\right)
  \;\bm{\Gamma}_h,
\end{equation}
where
\begin{itemize}[nosep]
  \item $\bm{\Gamma}_h = \sum_{i} \mathbf{r}_{t,i}\,\mathbf{r}_{t,i+h}^\top$
        is the $h$-th cross-autocovariance matrix of intraday returns.
  \item $k(\cdot)$: a kernel function satisfying $k(0) = 1$, $k(1) = 0$
        (e.g., Parzen kernel). Same options as the univariate case in
        Chapter~\ref{ch:noise}.
  \item $H$: bandwidth, controlling how many lags to include.
\end{itemize}
\end{definition}

\begin{keyresult}[PSD Guarantee Under Noise]
The multivariate realized kernel with a non-negative kernel function (e.g.,
Parzen) is guaranteed PSD by construction, even with non-synchronous trading
and microstructure noise \citep{BarndorffNielsenHansenLundeShephard2011}. This
is a major advantage over refresh-time RC and the Hayashi--Yoshida estimator.
The rate of convergence is $M^{-1/5}$ (slower than the noise-free $M^{-1/2}$),
the same price paid in the univariate case.
\end{keyresult}


% ──────────────────────────────────────────────────────────────
\section{DCC-GARCH}\label{sec:multivariate:dcc}
% ──────────────────────────────────────────────────────────────

\textbf{Where we are.}  You can now \emph{estimate} a daily covariance matrix
from high-frequency data; the next task is to \emph{forecast} tomorrow's.
DCC-GARCH is the multivariate analog of GARCH
(Chapter~\ref{ch:garch}): a parametric, easy-to-estimate baseline.

% ── TikZ: DCC decomposition ──
\begin{figure}[ht]
\centering
\begin{tikzpicture}[>=Stealth, node distance=1.5cm]
  % Blocks
  \node[flowblock, fill=defblue!10, minimum width=3cm] (vol) {
    \textbf{Step 1: Volatility}\\
    Univariate GARCH\\
    per asset $\to \sigma_{j,t}$
  };
  \node[flowblock, fill=keyorange!10, minimum width=3cm, right=5cm of vol] (corr) {
    \textbf{Step 2: Correlation}\\
    DCC dynamics\\
    $\to R_t$
  };
  \node[flowblock, fill=intgreen!10, minimum width=3cm, right=2.5cm of corr] (cov) {
    \textbf{Output}\\
    $\Sigma_t = D_t\, R_t\, D_t$
  };

  % Arrows
  \draw[arrow] (vol) -- node[above, font=\small, inner sep=2pt]{$D_t = \diag(\sigma_{1,t}, \ldots, \sigma_{p,t})$} (corr);
  \draw[arrow] (corr) -- (cov);
\end{tikzpicture}
\caption{DCC-GARCH decomposes the covariance matrix into volatilities (modeled
separately per asset) and a dynamic correlation matrix.}
\label{fig:dcc}
\end{figure}

\begin{keyidea}[DCC-GARCH \citep{Engle2002DCC}]
The Dynamic Conditional Correlation (DCC) model has two steps:

\textbf{Step 1 (volatilities).}  Fit a univariate GARCH(1,1) (or any variant)
to each asset's return series separately, producing conditional standard
deviations $\sigma_{j,t}$ for $j = 1, \ldots, p$. Stack them into a diagonal
matrix:
\[
  D_t = \diag(\sigma_{1,t}, \sigma_{2,t}, \ldots, \sigma_{p,t}).
\]

\textbf{Step 2 (correlations).}  Compute standardized residuals
$\mathbf{z}_t = D_t^{-1}\,\mathbf{r}_t$ and model their quasi-correlation:
\begin{equation}\label{eq:dcc}
  Q_t \;=\; (1 - \alpha - \beta)\,\bar{Q}
  \;+\; \alpha\,\mathbf{z}_{t-1}\,\mathbf{z}_{t-1}^\top
  \;+\; \beta\, Q_{t-1},
\end{equation}
where
\begin{itemize}[nosep]
  \item $\bar{Q}$: unconditional covariance of $\mathbf{z}_t$ (estimated from
        the full sample).
  \item $\alpha > 0$: weight on yesterday's innovation (analogous to ARCH
        coefficient).
  \item $\beta > 0$: persistence (analogous to GARCH coefficient).
  \item $\alpha + \beta < 1$: stationarity condition.
\end{itemize}
Then rescale to a proper correlation matrix:
\[
  R_t = \diag(Q_t)^{-1/2}\; Q_t\; \diag(Q_t)^{-1/2}.
\]
The conditional covariance matrix is $\Sigma_t = D_t\, R_t\, D_t$.
\end{keyidea}

\begin{intuition}[Why DCC is the ``HAR of Multivariate Vol'']
DCC is to multivariate volatility what HAR (Chapter~\ref{ch:har}) is to
univariate: not the most accurate model, but easy to estimate, hard to beat
consistently, and the first thing you should try. It scales well to large $p$
because Step~1 is embarrassingly parallel (one GARCH per asset) and Step~2 has
only two free parameters ($\alpha$, $\beta$) regardless of dimension.
\end{intuition}

\begin{warning}[DCC Limitations]
\begin{itemize}[nosep]
  \item Correlations follow a single $(\alpha, \beta)$ dynamic for all pairs.
        In reality, the IBM--MSFT correlation may move differently from the
        IBM--gold correlation.
  \item DCC uses daily returns, not intraday data. It ignores the richer
        information in realized covariance.
  \item The two-step estimation is not fully efficient (but it is consistent).
\end{itemize}
\end{warning}

\begin{projectconnection}[Why This Matters]
DCC-GARCH is the parametric baseline you should beat. It plays the same role
in multivariate vol forecasting that plain GARCH plays in univariate: simple,
well-understood, and surprisingly competitive. In a Diebold-Mariano test
comparing your ML covariance forecasts against DCC, a statistically significant
QLIKE improvement is the clearest evidence that high-frequency data and
nonlinear models add value beyond what a two-parameter correlation dynamic
can capture.
\end{projectconnection}


% ──────────────────────────────────────────────────────────────
\section{Wishart Autoregressive (WAR) Model}\label{sec:multivariate:war}
% ──────────────────────────────────────────────────────────────

\textbf{Where we are.}  DCC models correlations parametrically using daily
returns. Can you instead build a direct time-series model for the realized
covariance matrix, analogous to HAR for realized variance?

\begin{keyidea}[Wishart Autoregressive Model]
The WAR model treats the sequence of realized covariance matrices as a
matrix-variate time series. In its simplest form:
\begin{equation}\label{eq:war}
  \mathbf{RC}_{t+1} \;=\; C + A\,\mathbf{RC}_t\,A^\top + E_{t+1},
\end{equation}
where
\begin{itemize}[nosep]
  \item $C \in \R^{p \times p}$: intercept matrix (symmetric, PSD).
  \item $A \in \R^{p \times p}$: autoregressive coefficient matrix.
  \item $E_{t+1}$: innovation matrix (Wishart-distributed).
  \item The Wishart distribution is the matrix generalization of the
        chi-squared distribution and is the natural distribution for covariance
        matrices.
\end{itemize}
Higher-order and HAR-style variants (daily, weekly, monthly lags) follow
naturally:
\[
  \mathbf{RC}_{t+1} = C
  + A_d\,\mathbf{RC}_t\, A_d^\top
  + A_w\,\overline{\mathbf{RC}}_t^{(w)}\, A_w^\top
  + A_m\,\overline{\mathbf{RC}}_t^{(m)}\, A_m^\top
  + E_{t+1},
\]
where $\overline{\mathbf{RC}}_t^{(w)}$ and $\overline{\mathbf{RC}}_t^{(m)}$
are averages over the past 5 and 22 trading days.
\end{keyidea}

\begin{intuition}[In Plain English]
The WAR model is the matrix version of an AR(1) for time series. Instead of
saying ``tomorrow's variance equals a constant plus a fraction of today's
variance,'' it says ``tomorrow's covariance \emph{matrix} equals an intercept
matrix plus a transformation of today's covariance matrix.'' The sandwich
form $A\,\mathbf{RC}_t\,A^\top$ ensures the output remains symmetric (just as
$\mathbf{RC}_t$ is), and the Wishart distribution is the natural error
distribution for random matrices that must be positive semi-definite.
\end{intuition}

\begin{warning}[Parameter Explosion]
Each coefficient matrix $A$ has $p^2$ free parameters. For $p = 50$ assets,
$A$ alone has 2{,}500 parameters. With daily, weekly, and monthly lags, that
is 7{,}500 parameters plus the intercept. The sample sizes available
(typically 1{,}000--3{,}000 trading days) are far too small. WAR is practical
only for small $p$ (say, $p \leq 5$) unless you impose strong structure
(diagonal $A$, block-diagonal, etc.).
\end{warning}


% ──────────────────────────────────────────────────────────────
\section{HAR-DRD and Multivariate HARQ}\label{sec:multivariate:drd}
% ──────────────────────────────────────────────────────────────

\textbf{Where we are.}  WAR is theoretically clean but does not scale. The next
idea: decompose the covariance matrix into pieces that are easier to model
separately.

\begin{keyidea}[DRD Decomposition \citep{BollerslevPattonQuaedvlieg2018}]
Write the realized covariance matrix as
\begin{equation}\label{eq:drd}
  \mathbf{RC}_t \;=\; D_t\; R_t\; D_t,
\end{equation}
where
\begin{itemize}[nosep]
  \item $D_t = \diag\!\bigl(\sqrt{[\mathbf{RC}_t]_{11}}, \ldots,
        \sqrt{[\mathbf{RC}_t]_{pp}}\bigr)$: diagonal matrix of realized
        standard deviations.
  \item $R_t = D_t^{-1}\,\mathbf{RC}_t\, D_t^{-1}$: realized correlation
        matrix (unit diagonal).
\end{itemize}
Then model $D_t$ and $R_t$ separately:
\begin{itemize}[nosep]
  \item \textbf{Variances:} fit a univariate HAR (or HARQ, from
        Chapter~\ref{ch:har}) to each diagonal element $[\mathbf{RC}_t]_{jj}$.
  \item \textbf{Correlations:} fit a univariate HAR to each unique off-diagonal
        element of $R_t$, using Fisher $z$-transform to map $[-1,1]$ to
        $(-\infty, \infty)$.
\end{itemize}
Reassemble: $\widehat{\mathbf{RC}}_{t+1} = \hat{D}_{t+1}\,\hat{R}_{t+1}\,\hat{D}_{t+1}$.
\end{keyidea}

% ── TikZ: DRD decomposition schematic ──
\begin{figure}[ht]
\centering
\begin{tikzpicture}[>=Stealth, node distance=0.8cm,
    matbox/.style={draw, thick, minimum width=1.8cm, minimum height=1.8cm,
                   font=\small, align=center}]

  % RC matrix
  \node[matbox, fill=gray!15] (rc) {$\mathbf{RC}_t$\\{\scriptsize full}\\{\scriptsize covariance}};

  % Equals sign
  \node[right=0.6cm of rc, font=\Large] (eq) {$=$};

  % D matrix
  \node[matbox, right=0.6cm of eq, fill=defblue!20] (d1) {$D_t$\\{\scriptsize diagonal}\\{\scriptsize volatilities}};

  % Times sign
  \node[right=0.3cm of d1, font=\Large] (t1) {$\times$};

  % R matrix
  \node[matbox, right=0.3cm of t1, fill=keyorange!20] (r) {$R_t$\\{\scriptsize correlation}\\{\scriptsize matrix}};

  % Times sign
  \node[right=0.3cm of r, font=\Large] (t2) {$\times$};

  % D matrix again
  \node[matbox, right=0.3cm of t2, fill=defblue!20] (d2) {$D_t$\\{\scriptsize diagonal}\\{\scriptsize volatilities}};

  % Arrows to forecasting models below
  \node[flowblock, fill=defblue!10, below=1.2cm of d1, minimum width=2.2cm,
        font=\small, align=center] (har) {HAR / HARQ\\per asset};
  \node[flowblock, fill=keyorange!10, below=1.2cm of r, minimum width=2.2cm,
        font=\small, align=center] (harcorr) {HAR on Fisher\\$z$-transformed\\correlations};

  \draw[arrow] (d1) -- (har);
  \draw[arrow] (r) -- (harcorr);

  % Labels
  \node[defblue, font=\scriptsize, below=0.1cm of har] {$p$ univariate models};
  \node[keyorange, font=\scriptsize, below=0.1cm of harcorr]
    {$\binom{p}{2}$ univariate models};
\end{tikzpicture}
\caption{The DRD decomposition separates the covariance matrix into volatilities
(modeled individually with HAR/HARQ) and correlations (modeled with HAR on
Fisher $z$-transforms). Each piece is forecast with simple univariate models,
then reassembled.}
\label{fig:drd-decomp}
\end{figure}

\begin{keyresult}[Separate Modeling Beats Joint]
\citet{BollerslevPattonQuaedvlieg2018} show that the HARQ-DRD model, where
variances and correlations are modeled with separate HAR regressions and
measurement-error attenuation, significantly outperforms the direct vech-HAR
on all elements jointly. The HARQ variant (adding $\sqrt{\operatorname{RQ}}$
interactions from Chapter~\ref{ch:har}) further improves the variance
forecasts, especially on high-noise days.
\end{keyresult}

\begin{workedexample}{HAR-DRD for 2 Assets}
Suppose you have daily realized covariance matrices for stocks A and B.

\textbf{Step 1: extract components.}
\[
  \mathbf{RC}_t = \begin{pmatrix} 0.00040 & 0.00012 \\ 0.00012 & 0.00025 \end{pmatrix}.
\]
Then $D_t = \diag(0.0200, 0.0158)$ and
\[
  R_t = D_t^{-1}\,\mathbf{RC}_t\, D_t^{-1}
  = \begin{pmatrix} 1.000 & 0.379 \\ 0.379 & 1.000 \end{pmatrix}.
\]

\textbf{Step 2: model correlation.}  Apply Fisher $z$-transform:
$z_t = \frac{1}{2}\ln\!\bigl(\frac{1 + 0.379}{1 - 0.379}\bigr) \approx 0.399$.
Fit HAR to the $\{z_t\}$ series. Invert the transform on the forecast to get
$\hat{\rho}_{t+1}$.
\end{workedexample}

\begin{warning}[PSD Not Guaranteed]
The separately forecasted $\hat{R}_{t+1}$ is not guaranteed to be a valid
correlation matrix (it may have eigenvalues outside $[0,1]$ or a diagonal
different from 1). In practice, you project onto the nearest correlation
matrix, e.g., via the alternating projections algorithm of Higham (2002).
\end{warning}


% ──────────────────────────────────────────────────────────────
\section{Cholesky-HAR}\label{sec:multivariate:cholesky}
% ──────────────────────────────────────────────────────────────

\textbf{Where we are.}  HAR-DRD decomposes the matrix into variances and
correlations. Cholesky-HAR takes a different decomposition that guarantees PSD
by construction.

\begin{keyidea}[Cholesky-HAR \citep{ChiriacVoev2011}]
Every PSD matrix $\mathbf{RC}_t$ has a unique Cholesky decomposition:
\[
  \mathbf{RC}_t = L_t\, L_t^\top,
\]
where $L_t$ is lower triangular with positive diagonal entries. The idea:
\begin{enumerate}[nosep]
  \item Compute $L_t$ for each day.
  \item Take the log of diagonal elements (to ensure positivity upon
        exponentiation).
  \item Stack all $p(p+1)/2$ unique elements of the modified $L_t$ into a
        vector.
  \item Fit a separate HAR model to each element.
  \item Forecast, exponentiate diagonals, unstack into $\hat{L}_{t+1}$,
        and reassemble:
        $\widehat{\mathbf{RC}}_{t+1} = \hat{L}_{t+1}\,\hat{L}_{t+1}^\top$.
\end{enumerate}
Since any lower triangular matrix with positive diagonal gives a valid PSD
matrix via $LL^\top$, the forecast is PSD by construction.
\end{keyidea}

\begin{workedexample}{Cholesky-HAR for 2 Assets}
Using the same $\mathbf{RC}_t$ from the DRD example:
\[
  \mathbf{RC}_t = \begin{pmatrix} 0.00040 & 0.00012 \\ 0.00012 & 0.00025 \end{pmatrix}.
\]

\textbf{Step 1: Cholesky decomposition.}
\[
  L_t = \begin{pmatrix} 0.02000 & 0 \\ 0.00600 & 0.01463 \end{pmatrix}.
\]
Check: $l_{11} = \sqrt{0.00040} = 0.02$,\;
$l_{21} = 0.00012 / 0.02 = 0.006$,\;
$l_{22} = \sqrt{0.00025 - 0.006^2} = \sqrt{0.000214} \approx 0.01463$.

\textbf{Step 2: transform.}
Log-transform diagonal: $(\ln 0.02, \ln 0.01463) \approx (-3.912, -4.225)$.
Off-diagonal left as-is: $0.006$.

\textbf{Step 3: fit HAR to each of the 3 elements.}

\textbf{Step 4: forecast and reassemble.}
Suppose the HAR forecasts, after exponentiating the diagonals, yield
$\hat{l}_{11} \approx 0.02020$, $\hat{l}_{21} = 0.00580$, $\hat{l}_{22} \approx 0.01483$.
Then
\[
  \widehat{\mathbf{RC}}_{t+1}
  = \begin{pmatrix} 0.02020 & 0 \\ 0.00580 & 0.01483 \end{pmatrix}
    \begin{pmatrix} 0.02020 & 0.00580 \\ 0 & 0.01483 \end{pmatrix}
  = \begin{pmatrix} 0.000408 & 0.000117 \\ 0.000117 & 0.000254 \end{pmatrix}.
\]
This is PSD by construction (it is $LL^\top$).
\end{workedexample}


% ──────────────────────────────────────────────────────────────
\section{Graph-Based Methods}\label{sec:multivariate:graph}
% ──────────────────────────────────────────────────────────────

\textbf{Where we are.}  All methods so far treat each element (or factor) of
the covariance matrix independently. Graph-based methods exploit the fact that
assets are \emph{connected}: if asset A's volatility spills over to asset B
(Chapter~\ref{ch:spillovers} formalizes this), modeling them jointly through
a graph should help.

% ── TikZ: Asset graph ──
\begin{figure}[ht]
\centering
\begin{tikzpicture}[>=Stealth]
  % Nodes
  \node[obs, fill=defblue!20] (A) at (0, 2) {AAPL};
  \node[obs, fill=defblue!20] (B) at (3, 3.2) {MSFT};
  \node[obs, fill=defblue!20] (C) at (5.5, 1.5) {GOOG};
  \node[obs, fill=warnred!20] (D) at (1.5, -0.5) {XOM};
  \node[obs, fill=warnred!20] (E) at (4.5, -0.8) {CVX};

  % Edges with weights
  \draw[thick, defblue] (A) -- node[above, font=\scriptsize, sloped]{$w=0.72$} (B);
  \draw[thick, defblue] (B) -- node[above, font=\scriptsize, sloped]{$w=0.68$} (C);
  \draw[thick, defblue!40] (A) -- node[below, font=\scriptsize, sloped]{$w=0.35$} (C);
  \draw[thick, warnred] (D) -- node[below, font=\scriptsize, sloped]{$w=0.81$} (E);
  \draw[thick, gray!60] (A) -- node[left, font=\scriptsize]{$w=0.15$} (D);
  \draw[thick, gray!60] (C) -- node[right, font=\scriptsize]{$w=0.18$} (E);
  \draw[thick, gray!60] (B) -- node[right, font=\scriptsize]{$w=0.12$} (E);

  % Legend
  \node[right, font=\small] at (6.5, 3) {Edge weight = realized};
  \node[right, font=\small] at (6.5, 2.5) {correlation $|\rho_{ij}|$};
  \node[right, font=\small, defblue] at (6.5, 1.7) {Tech cluster};
  \node[right, font=\small, warnred] at (6.5, 1.2) {Energy cluster};
\end{tikzpicture}
\caption{Assets as nodes in a graph with edge weights derived from realized
correlations. Within-sector edges (tech, energy) are strong; cross-sector edges
are weak. Graph-based models exploit this structure.}
\label{fig:asset-graph}
\end{figure}

\begin{keyidea}[Graph-HAR \citep{ZhangPuCucuringuDong2024}]
Graph-HAR augments the standard HAR model with graph-based features. The
adjacency matrix $W$ is constructed from realized correlations (or partial
correlations). For each asset $j$:
\begin{equation}\label{eq:graph-har}
  \RV_{j,t+1} \;=\; \beta_0 + \beta_d\,\RV_{j,t}
  + \beta_w\,\RV_{j,t}^{(w)}
  + \beta_m\,\RV_{j,t}^{(m)}
  + \gamma\, \sum_{k \neq j} W_{jk}\,\RV_{k,t}
  + \varepsilon_{j,t+1},
\end{equation}
where
\begin{itemize}[nosep]
  \item The first three terms are the standard HAR (Chapter~\ref{ch:har}).
  \item $\sum_{k \neq j} W_{jk}\,\RV_{k,t}$: a graph-weighted average of
        neighbors' volatilities. This is exactly one step of graph diffusion.
  \item $\gamma$: spillover coefficient. If $\gamma > 0$, high volatility in
        neighbors predicts higher volatility for asset $j$.
  \item $W$: adjacency matrix, typically from thresholded correlation or LASSO
        partial correlation.
\end{itemize}
\end{keyidea}

\begin{keyidea}[GNN for Covariance \citep{ZhangCucuringuDong2023}]
Going further, Graph Neural Networks (GNNs) replace the linear spillover term
with learnable message-passing layers:
\begin{enumerate}[nosep]
  \item \textbf{Node features:} each asset's HAR-style inputs (daily, weekly,
        monthly RV, plus optional firm characteristics).
  \item \textbf{Graph structure:} adjacency from realized correlation, GICS
        sector membership, or learned adaptively.
  \item \textbf{Message passing:} each node aggregates features from its
        neighbors through learned weight matrices (as in Chapter~\ref{ch:deeplearning}).
  \item \textbf{Output:} node-level volatility forecasts or, with a symmetric
        readout layer, full covariance matrix forecasts.
\end{enumerate}
The advantage over Graph-HAR: nonlinear aggregation and multi-hop
information flow (asset A learns from asset C through intermediary B).
\end{keyidea}

% ── TikZ: GNN message passing ──
\begin{figure}[ht]
\centering
\begin{tikzpicture}[>=Stealth, node distance=2.2cm,
    assetnode/.style={circle, draw, thick, minimum size=1.1cm, font=\small},
    msgbox/.style={rectangle, draw, rounded corners, thick, font=\scriptsize,
                   fill=keyorange!10, minimum width=1.4cm, align=center}]

  % Layer 0 (input)
  \node[assetnode, fill=defblue!20] (a0) {A};
  \node[assetnode, fill=defblue!20, below left=1.5cm and 2cm of a0] (b0) {B};
  \node[assetnode, fill=defblue!20, below right=1.5cm and 2cm of a0] (c0) {C};

  % Edges
  \draw[thick, gray] (a0) -- (b0);
  \draw[thick, gray] (a0) -- (c0);
  \draw[thick, gray!40, dashed] (b0) -- (c0);

  % Layer label
  \node[above=0.3cm of a0, font=\small\bfseries] {Input features};
  \node[below=0.1cm of a0, font=\scriptsize, defblue] {$\RV_A^{(d,w,m)}$};
  \node[below=0.1cm of b0, font=\scriptsize, defblue] {$\RV_B^{(d,w,m)}$};
  \node[below=0.1cm of c0, font=\scriptsize, defblue] {$\RV_C^{(d,w,m)}$};

  % Arrow to aggregation
  \node[right=3.5cm of a0] (mid) {};
  \node[msgbox, above=0.2cm of mid] (agg) {Aggregate\\neighbors};
  \node[msgbox, below=0.8cm of agg] (upd) {Update\\node $A$};

  \draw[arrow, defblue] (b0.east) .. controls +(1.5,0.5) .. (agg.west);
  \draw[arrow, defblue] (c0.north east) .. controls +(1,1) .. (agg.south west);
  \draw[arrow] (agg) -- (upd);
  \draw[arrow, dashed, gray] (a0) -- (upd);

  % Output
  \node[assetnode, fill=intgreen!20, right=1.5cm of upd] (out)
    {$\hat{h}_A$};
  \draw[arrow] (upd) -- (out);
  \node[below=0.2cm of out, font=\scriptsize, intgreen, align=center]
    {Updated\\representation};

  % Annotation
  \node[right=0.3cm of out, font=\scriptsize, align=left] (ann)
    {After $K$ layers:\\asset $A$ has\\``seen'' $K$-hop\\neighbors};
\end{tikzpicture}
\caption{GNN message passing for one node. Asset A aggregates volatility
features from its graph neighbors (B and C), then updates its own
representation. After $K$ message-passing layers, each node has information
from all assets within $K$ hops, capturing multi-step spillover effects.}
\label{fig:gnn-message}
\end{figure}


% ──────────────────────────────────────────────────────────────
\section{CNN-RCOV}\label{sec:multivariate:cnn}
% ──────────────────────────────────────────────────────────────

\textbf{Where we are.}  Graph methods treat assets as nodes. An alternative:
treat the realized covariance matrix itself as an image.

\begin{keyidea}[Covariance Matrices as Images]
A $p \times p$ realized covariance matrix is a symmetric, real-valued ``image.''
If you sort assets by sector (tech, energy, financials, \ldots), block structure
emerges: within-sector blocks have high covariance, cross-sector blocks have
low covariance. Convolutional neural networks (CNNs, Chapter~\ref{ch:deeplearning})
are designed to detect such local spatial patterns.

\textbf{Architecture.}  Stack $T$ past covariance matrices as ``channels''
(like RGB channels in image classification) and use 2D convolutions to
extract temporal and cross-sectional patterns:
\begin{enumerate}[nosep]
  \item Input: $\mathbf{RC}_{t}, \mathbf{RC}_{t-1}, \ldots, \mathbf{RC}_{t-T+1}$
        $\in \R^{T \times p \times p}$.
  \item 2D convolutional layers detect local block patterns.
  \item Fully connected layers map to the $p(p+1)/2$ unique elements of
        $\widehat{\mathbf{RC}}_{t+1}$.
\end{enumerate}
\end{keyidea}

\begin{warning}[CNN-RCOV Caveats]
\begin{itemize}[nosep]
  \item The literature is thin; this is more ``conceptually appealing'' than
        empirically proven at scale.
  \item Ordering assets by sector is a design choice. Different orderings yield
        different spatial patterns, and CNNs are not permutation-invariant.
        GNNs (Section~\ref{sec:multivariate:graph}) handle this more
        naturally.
  \item The output is not guaranteed PSD; post-hoc projection is required.
  \item For large $p$, the output layer has $O(p^2)$ parameters, which grows
        fast.
\end{itemize}
\end{warning}


% ──────────────────────────────────────────────────────────────
\section{Geometric Deep Learning on the SPD Manifold}\label{sec:multivariate:spd}
% ──────────────────────────────────────────────────────────────

\textbf{Where we are.}  Every method so far either ignores the PSD constraint
or enforces it with post-hoc projection. SPDNet takes a fundamentally different
approach: work directly on the manifold where covariance matrices live.

\begin{prereq}[What is a Manifold?]
A manifold is a curved space that locally looks like Euclidean space, much like
the surface of a sphere locally looks flat. The set of $p \times p$ symmetric
positive definite matrices (denoted $\mathcal{S}_{++}^p$) forms a smooth
manifold. It is \emph{not} a flat (Euclidean) space: the straight line
between two PSD matrices can pass through matrices with zero or negative
eigenvalues, leaving the PSD cone.
\end{prereq}

% ── TikZ: SPD manifold ──
\begin{figure}[ht]
\centering
\begin{tikzpicture}[>=Stealth, scale=1.2]
  % PSD cone (simplified 2D projection)
  \fill[intgreen!10] (0.5,0) .. controls (1,3) and (4,3.5) .. (6,3)
    .. controls (5,1) and (2,0.5) .. (0.5,0);
  \draw[intgreen, thick] (0.5,0) .. controls (1,3) and (4,3.5) .. (6,3)
    .. controls (5,1) and (2,0.5) .. (0.5,0);
  \node[intgreen, font=\small] at (3.5, 2.8) {$\mathcal{S}_{++}^p$ (PSD cone)};

  % Two PSD matrices
  \fill[defblue] (1.5, 1.5) circle (3pt) node[above left, font=\small]{$\Sigma_A$};
  \fill[defblue] (4.5, 1.8) circle (3pt) node[above right, font=\small]{$\Sigma_B$};

  % Euclidean line (leaves cone)
  \draw[warnred, thick, dashed] (1.5, 1.5) -- (4.5, 1.8);
  \node[warnred, font=\scriptsize, below] at (3.0, 1.2) {Euclidean interpolation};
  \fill[warnred] (3.0, 1.65) circle (2pt);
  \node[warnred, font=\scriptsize] at (3.0, 0.8) {(may leave PSD cone!)};

  % Geodesic (stays inside)
  \draw[defblue, thick] (1.5, 1.5) .. controls (2.5, 2.5) and (3.5, 2.6) .. (4.5, 1.8);
  \node[defblue, font=\scriptsize] at (2.8, 2.6) {Riemannian geodesic};
  \node[defblue, font=\scriptsize] at (2.8, 2.3) {(stays inside PSD cone)};

  % Boundary
  \draw[gray, thick, ->] (0, -0.2) -- (6.5, -0.2);
  \node[gray, font=\scriptsize, right] at (6.5, -0.2) {boundary: $\det = 0$};
\end{tikzpicture}
\caption{The SPD manifold. A straight (Euclidean) line between two PSD matrices
$\Sigma_A$ and $\Sigma_B$ can pass through singular matrices (red dashed). The
Riemannian geodesic (blue curve) stays inside the PSD cone by definition.}
\label{fig:spd-manifold}
\end{figure}

\begin{definition}[Affine-Invariant Riemannian Metric on $\mathcal{S}_{++}^p$]
The standard Riemannian metric on $\mathcal{S}_{++}^p$ is
\[
  d(\Sigma_A, \Sigma_B)
  \;=\; \bigl\|\log\bigl(\Sigma_A^{-1/2}\,\Sigma_B\,\Sigma_A^{-1/2}\bigr)\bigr\|_F,
\]
where $\log$ is the matrix logarithm and $\|\cdot\|_F$ is the Frobenius norm.
This distance is affine-invariant: $d(A\Sigma_A A^\top, A\Sigma_B A^\top) =
d(\Sigma_A, \Sigma_B)$ for any invertible $A$.
\end{definition}

\begin{keyidea}[SPDNet]
SPDNet replaces standard neural network operations with Riemannian analogs that
preserve positive definiteness:
\begin{enumerate}[nosep]
  \item \textbf{BiMap layer} (analog of linear layer):
        $\Sigma \mapsto W\,\Sigma\, W^\top$, where $W$ is a learnable
        $d \times p$ matrix with $d \leq p$. If $\Sigma \succ 0$ and $W$ has
        full row rank, then $W\Sigma W^\top \succ 0$. This also reduces
        dimension from $p$ to $d$.
  \item \textbf{ReEig layer} (analog of ReLU):
        eigendecompose $\Sigma = U\Lambda U^\top$, then
        $\Sigma \mapsto U\,\max(\Lambda, \epsilon I)\, U^\top$. Clips small
        eigenvalues to $\epsilon > 0$, keeping the matrix strictly PSD.
  \item \textbf{LogEig layer} (analog of final feature extraction):
        $\Sigma \mapsto U\,\log(\Lambda)\,U^\top$. Maps from the SPD manifold
        to the tangent space (symmetric matrices), where standard Euclidean
        operations apply.
\end{enumerate}
The output of LogEig lives in a flat space, so you can apply a standard linear
classifier or regressor on top.
\end{keyidea}

\begin{keyresult}[SPDNet: PSD by Design]
Because every layer in SPDNet maps PSD inputs to PSD outputs, intermediate
representations are always valid covariance matrices. No post-hoc projection
is needed. This is the only deep learning architecture for covariance
forecasting with a built-in PSD guarantee.
\end{keyresult}


% ──────────────────────────────────────────────────────────────
\section{The PSD Constraint}\label{sec:multivariate:psd}
% ──────────────────────────────────────────────────────────────

\textbf{Where we are.}  Every multivariate volatility model must produce a
valid (positive semi-definite) covariance matrix. Some methods guarantee this;
others require fixing after the fact. This section provides a systematic
comparison.

\begin{intuition}[Why PSD Matters]
A covariance matrix that is not PSD implies negative variance for some portfolio.
Specifically, if $\Sigma$ has a negative eigenvalue with eigenvector
$\mathbf{v}$, then the portfolio $\mathbf{v}$ has predicted variance
$\mathbf{v}^\top \Sigma\, \mathbf{v} < 0$. This is nonsensical and causes
optimizers to blow up (infinite leverage on the ``negative variance'' direction).
\end{intuition}

\begin{table}[ht]
\centering
\caption{PSD guarantees across multivariate volatility methods.}
\label{tab:psd}
\begin{tabular}{lll}
\toprule
\textbf{Method} & \textbf{PSD Guarantee} & \textbf{Mechanism} \\
\midrule
DCC-GARCH (Sec.~\ref{sec:multivariate:dcc})
  & Yes & By construction ($D_t R_t D_t$, $R_t$ from rescaling) \\
WAR (Sec.~\ref{sec:multivariate:war})
  & No  & Post-hoc projection needed \\
HAR-DRD (Sec.~\ref{sec:multivariate:drd})
  & No  & Post-hoc projection on $\hat{R}_{t+1}$ \\
Cholesky-HAR (Sec.~\ref{sec:multivariate:cholesky})
  & Yes & $\hat{L}\hat{L}^\top$ is PSD by construction \\
Graph-HAR (Sec.~\ref{sec:multivariate:graph})
  & No  & Post-hoc projection needed \\
CNN-RCOV (Sec.~\ref{sec:multivariate:cnn})
  & No  & Post-hoc projection needed \\
SPDNet (Sec.~\ref{sec:multivariate:spd})
  & Yes & Riemannian operations preserve PSD \\
\bottomrule
\end{tabular}
\end{table}

\begin{keyidea}[Post-Hoc PSD Projection]
When a method does not guarantee PSD, the standard fix is eigenvalue clipping:
\begin{enumerate}[nosep]
  \item Eigendecompose: $\hat{\Sigma} = U\Lambda U^\top$.
  \item Set $\tilde{\Lambda} = \max(\Lambda, 0)$ (zero out negative
        eigenvalues).
  \item Reconstruct: $\tilde{\Sigma} = U\tilde{\Lambda} U^\top$.
\end{enumerate}
This gives the nearest PSD matrix in Frobenius norm. The cost: it introduces
bias, and the projected matrix no longer minimizes whatever loss function the
model was trained on. For correlation matrices, the Higham (2002) alternating
projection algorithm additionally enforces unit diagonal.
\end{keyidea}

\begin{warning}[PSD Violations in Practice]
For small $p$ (2--5 assets), PSD violations are rare and small. For large $p$
(50--500), they are common and can be severe, especially with element-wise
models (HAR-DRD, Graph-HAR) that do not enforce cross-element consistency.
If PSD is critical for your application (portfolio optimization, risk
management), prefer methods with built-in guarantees: DCC, Cholesky-HAR, or
SPDNet.
\end{warning}


% ──────────────────────────────────────────────────────────────
\section{Summary}\label{sec:multivariate:summary}
% ──────────────────────────────────────────────────────────────

Key takeaways from this chapter:

\begin{itemize}[nosep]
  \item Realized covariance (RC) extends realized variance to matrices via
        outer products of intraday return vectors (\eqref{eq:rc}).
  \item Non-synchronous trading biases RC toward zero (Epps effect).
        Refresh-time sampling and the Hayashi--Yoshida estimator
        (\eqref{eq:hy}) address this, with HY being more data-efficient.
  \item The multivariate realized kernel (\eqref{eq:mrk}) handles both noise
        and non-synchronicity while guaranteeing PSD
        \citep{BarndorffNielsenHansenLundeShephard2011}.
  \item DCC-GARCH \citep{Engle2002DCC} is the standard parametric baseline:
        easy to estimate, scales to large $p$, but uses only daily data and
        imposes a single correlation dynamic.
  \item WAR models RC directly as a matrix time series, but suffers from
        parameter explosion for $p > 5$.
  \item HAR-DRD \citep{BollerslevPattonQuaedvlieg2018} decomposes RC into
        variances and correlations, models each with HAR, and outperforms
        both joint modeling and DCC.
  \item Cholesky-HAR \citep{ChiriacVoev2011} models Cholesky factors with HAR,
        guaranteeing PSD by construction.
  \item Graph-HAR \citep{ZhangPuCucuringuDong2024} and GNN methods
        \citep{ZhangCucuringuDong2023} model volatility spillovers through
        asset graphs, capturing network effects that element-wise models miss.
  \item CNN-RCOV treats RC matrices as images; conceptually appealing but
        empirically limited and not permutation-invariant.
  \item SPDNet operates on the Riemannian manifold of PSD matrices, the only
        deep learning approach with a built-in PSD guarantee.
  \item Methods either guarantee PSD (DCC, Cholesky-HAR, SPDNet) or require
        post-hoc eigenvalue projection, which introduces bias.
  \item For most applications, start with HAR-DRD (best accuracy in large-scale
        comparisons) or Cholesky-HAR (if PSD is essential). Use DCC-GARCH as
        the parametric baseline. Graph and manifold methods are promising but
        still maturing.
\end{itemize}

\begin{table}[ht]
\centering
\caption{Key results: multivariate volatility methods.}
\label{tab:multivariate-summary}
\begin{tabular}{p{3.2cm} p{3cm} p{3cm} p{3cm}}
\toprule
\textbf{Method} & \textbf{Data Input} & \textbf{Key Advantage} & \textbf{Key Limitation} \\
\midrule
DCC-GARCH & Daily returns & Scales to large $p$; PSD guaranteed & Single correlation dynamic; ignores HF data \\
WAR & Realized covariance & Theoretically clean matrix AR & $O(p^2)$ parameters \\
HAR-DRD & Realized covariance & Best accuracy (BPQ 2018); flexible & PSD not guaranteed \\
Cholesky-HAR & Realized covariance & PSD by construction & Cholesky ordering arbitrary \\
Graph-HAR / GNN & RV + graph & Models spillovers & Graph construction is a design choice \\
CNN-RCOV & RC as image & Detects block structure & Not permutation-invariant \\
SPDNet & RC on manifold & PSD by design; principled geometry & Complex; limited empirical evidence \\
\bottomrule
\end{tabular}
\end{table}

\bigskip
\noindent\textbf{Next.}  Chapter~\ref{ch:spillovers} formalizes the volatility
spillovers that Graph-HAR exploits, using the Diebold--Yilmaz connectedness
framework and network topology.

